\begin{onehalfspace}

\section{Il Limite del Calcolo Classico}

Per oltre cinquant'anni, lo sviluppo dell'informatica è stato guidato da una legge empirica enunciata nel 1965 da Gordon Moore\footnote{\textbf{Gordon Earle Moore} (1929--2023), cofondatore di Intel, osservò nel 1965 che il numero di transistor integrabili su un chip raddoppiava approssimativamente ogni due anni, con un conseguente aumento esponenziale delle prestazioni computazionali. La Legge di Moore non è una legge fisica, ma una previsione industriale che ha orientato decenni di investimenti tecnologici.}: il numero di transistor per chip raddoppia ogni due anni, e con esso la potenza di calcolo disponibile. Questa crescita esponenziale ha permesso di trasformare calcolatori grandi quanto un'intera stanza nei dispositivi tascabili di oggi. Tuttavia, la Legge di Moore si sta scontrando con i limiti fisici fondamentali della materia: i transistor moderni hanno dimensioni di pochi nanometri, avvicinandosi alla scala atomica oltre la quale gli effetti quantistici -- lo stesso fenomeno che si vorrebbe evitare -- diventano incontrollabili e inevitabili.

Parallelamente al problema fisico dell'hardware, esiste una barriera di natura puramente matematica: esistono classi di problemi computazionali per i quali nessun algoritmo classico è in grado di fornire una soluzione in tempi ragionevoli, \textit{indipendentemente} dalla potenza di calcolo disponibile. Questi problemi -- fattorizzazione di grandi interi, ricerca in spazi combinatoriali enormi, simulazione di sistemi molecolari complessi -- resistono a qualunque approccio classico non perché i computer non siano abbastanza veloci, ma perché la complessità intrinseca del problema cresce più rapidamente di qualunque risorsa classica concepibile.

È in questo contesto che nasce il \textbf{calcolo quantistico}: la visione radicale che i principi della meccanica quantistica -- \textit{sovrapposizione}, \textit{entanglement} e \textit{interferenza} -- possano essere sfruttati non come ostacoli da evitare, ma come risorse computazionali di straordinaria potenza. L'intuizione originale fu di Richard Feynman\footnote{\textbf{Richard Phillips Feynman} (1918--1988), fisico teorico statunitense, Premio Nobel per la Fisica nel 1965. Nel suo celebre articolo del 1982 \textit{``Simulating Physics with Computers''} \cite{feynman1982}, Feynman argomentò che un computer classico non può simulare efficientemente un sistema quantistico senza un overhead esponenziale, e propose che un ``computer quantistico'' -- uno che segua le leggi della meccanica quantistica -- potrebbe farlo in modo naturalmente efficiente.}, che nel 1982 osservò come la simulazione di sistemi quantistici richiedesse risorse esponenziali su macchine classiche, e propose che un computer basato sugli stessi principi fisici potesse simulare la natura stessa in modo efficiente \cite{feynman1982}. David Deutsch formalizzò nel 1985 il primo modello teorico rigoroso di computer quantistico universale \cite{deutsch1985}, aprendo la strada a un nuovo paradigma computazionale.

\section{La Rivoluzione Algoritmica degli Anni Novanta}

Negli anni successivi alla proposta di Deutsch, la comunità scientifica si interrogò sulla effettiva esistenza di problemi che un computer quantistico potesse risolvere in modo \textit{probabilmente} più efficiente di qualunque macchina classica. La risposta arrivò in forma dirompente nel corso di pochi anni.

Nel 1992, Daniel Simon\footnote{L'\textbf{algoritmo di Simon} (1994, pubblicato nel 1997) fu il primo a dimostrare uno speedup \textit{esponenziale} rispetto a qualunque algoritmo classico probabilistico, stabilendo la prova di principio che i computer quantistici potessero essere genuinamente più potenti. Shor riconobbe che la sua tecnica si ispirava direttamente alle idee di Simon.} identificò un problema artificiale per il quale un algoritmo quantistico era esponenzialmente più veloce di qualunque algoritmo classico. Questo risultato rimase per alcuni anni una curiosità teorica, fino a quando, nel \textbf{1994}, Peter Shor pubblicò un risultato che scosse le fondamenta della sicurezza informatica mondiale: un algoritmo quantistico in grado di fattorizzare un intero di $n$ bit in tempo \textit{polinomiale} $O(n^3)$ \cite{shor1994}, a fronte della complessità sub-esponenziale ma super-polinomiale dei migliori algoritmi classici noti. Le implicazioni furono immediate e devastanti: l'intero edificio della crittografia a chiave pubblica -- \ac{RSA}, Diffie-Hellman, le firme digitali \ac{ECDSA} -- si fondava proprio sull'assunzione che la fattorizzazione fosse computazionalmente intrattabile. Un computer quantistico sufficientemente grande avrebbe reso obsoleta, in un solo colpo, l'intera infrastruttura crittografica su cui poggia Internet.

Due anni dopo, nel \textbf{1996}, Lov Grover presentò un algoritmo di tutt'altra natura, ma di eguale importanza \cite{grover1996}. Mentre Shor aveva risolto un problema specifico della teoria dei numeri, Grover affrontò una sfida universale: la \textbf{ricerca non strutturata}. Dato un insieme di $N$ elementi senza alcuna struttura o ordinamento, trovare quello che soddisfa una certa proprietà richiede, nel caso classico, un numero di interrogazioni proporzionale a $N$. Grover dimostrò che un algoritmo quantistico può fare lo stesso lavoro in sole $O(\sqrt{N})$ interrogazioni -- uno \textit{speedup quadratico} -- e che questo risultato è \textit{ottimale}: nessun algoritmo quantistico può fare di meglio. Sebbene meno spettacolare dello speedup esponenziale di Shor, il vantaggio di Grover è universale, applicabile a qualunque problema di ricerca, e porta con sé implicazioni profonde per la crittografia simmetrica, l'ottimizzazione combinatoria e la ricerca in grandi basi di dati.

Questi due algoritmi -- Shor e Grover -- rappresentano ancora oggi i pilastri dell'informatica algoritmica quantistica: il primo come dimostrazione del potere \textit{esponenziale} del calcolo quantistico su problemi strutturati, il secondo come strumento \textit{universale} di accelerazione quadratica applicabile al problema computazionale più fondamentale.

\section{Lo Stato dell'Arte: dall'Era NISQ ai Computer Fault-Tolerant}

Nonostante la solidità teorica di questi risultati, la realizzazione pratica di computer quantistici di scala sufficiente ad eseguirli su istanze realmente significative incontra ostacoli formidabili. I sistemi quantistici fisici sono intrinsecamente fragili: qualunque interazione non controllata con l'ambiente esterno -- vibrazioni termiche, campi elettromagnetici, imperfezioni nei materiali -- provoca fenomeni di \textbf{decoerenza} -- ossia la perdita progressiva della coerenza quantistica causata dall'entanglement indesiderato tra il qubit e i gradi di libertà ambientali, che sopprime le interferenze quantistiche riducendo il sistema a una miscela statistica classica -- degradando lo stato quantistico in modo irreversibile prima che il calcolo possa completarsi. A differenza di un bit classico, che può essere protetto dalla ridondanza, un qubit non può essere copiato (teorema del no-cloning \cite{wootters1982}) né letto senza distruggerlo: questo rende la correzione degli errori quantistici un problema radicalmente diverso e molto più complesso di quella classica.

Il panorama tecnologico attuale è dominato da un insieme variegato di piattaforme hardware in competizione: \textbf{qubit superconduttori} (\ac{IBM}, Google, Rigetti), \textbf{ioni intrappolati} (IonQ, Quantinuum), \textbf{fotoni} (PsiQuantum, Xanadu) e \textbf{spin in silicio} (Intel). Ciascuna architettura presenta vantaggi e svantaggi in termini di tempi di coerenza, fedeltà delle porte, scalabilità e temperatura operativa.

Un momento simbolico di questa corsa fu raggiunto nel \textbf{2019}, quando il team di Google rivendicò il raggiungimento della cosiddetta \textit{quantum supremacy}\footnote{Il termine \textit{quantum supremacy}, coniato da John Preskill nel 2012, indica il punto in cui un computer quantistico esegue un calcolo specifico in un tempo che nessun supercomputer classico potrebbe eguagliare in modo pratico. La rivendicazione di Google con il processore Sycamore è rimasta controversa: IBM sostenne che il calcolo avrebbe potuto essere eseguito classicamente in 2.5 giorni anziché 10.000 anni, con un algoritmo migliorato. Questo dibattito evidenzia la difficoltà di definire e misurare il vantaggio quantistico.} con il processore Sycamore a 53 qubit \cite{arute2019}: un calcolo specifico fu completato in 200 secondi, mentre le stime attribuivano al più potente supercomputer classico un tempo di circa 10.000 anni per lo stesso compito. Nel 2023, IBM ha presentato il processore \textit{Condor} a 1.121 qubit e ha delineato una roadmap verso sistemi con oltre 100.000 qubit entro la fine del decennio.

Tuttavia, la distanza tra i dispositivi attuali e un computer quantistico \textit{fault-tolerant} -- in grado di eseguire Shor su istanze RSA reali -- è ancora enorme. I processori odierni operano nel regime cosiddetto \ac{NISQ}, un termine coniato da John Preskill nel 2018 \cite{preskill2018} per descrivere macchine con decine o centinaia di qubit fisici, affette da errori significativi che limitano la profondità dei circuiti eseguibili. In questo scenario, la ricerca si concentra su due direzioni parallele e complementari:

\begin{itemize}
    \item \textbf{\ac{QEC}}: tecniche per codificare qubit logici in qubit fisici ridondanti, in modo da rilevare e correggere gli errori durante il calcolo, a costo di un overhead fisico significativo (centinaia o migliaia di qubit fisici per qubit logico);

    \item \textbf{Noise mitigation}: tecniche che, senza richiedere overhead hardware, riducono l'effetto del rumore a livello di post-processing classico, estendendo la portata dei circuiti eseguibili su hardware NISQ.
\end{itemize}

\clearpage

\section{Motivazione e Contributo del Presente Lavoro}

In questo contesto scientifico e tecnologico si inserisce il presente lavoro di tesi. L'algoritmo di Shor rappresenta uno dei risultati più significativi e dirompenti dell'informatica quantistica: la sua capacità di fattorizzare interi in tempo polinomiale minaccia direttamente le fondamenta della crittografia moderna. Tuttavia, la sua esecuzione su hardware NISQ è ostacolata da requisiti stringenti in termini di profondità del circuito, numero di qubit e fedeltà delle porte -- rendendo il rumore quantistico il principale ostacolo alla sua realizzazione pratica.

Questo lavoro si propone di affrontare tale sfida integrando due ambiti in rapida evoluzione: il \textbf{calcolo quantistico} e i \textbf{modelli di intelligenza artificiale}. L'obiettivo centrale è implementare l'algoritmo di Shor su simulatori quantistici realistici, analizzarne sistematicamente la vulnerabilità al rumore nelle sue componenti critiche -- Quantum Fourier Transform, phase estimation e modular exponentiation -- e sviluppare strategie di ottimizzazione e mitigazione basate su modelli di machine learning.

Il contributo originale consiste nella sperimentazione diretta con il framework Qiskit, nella costruzione e analisi di circuiti per l'algoritmo di Shor in presenza di rumore realistico, e nell'applicazione di tecniche di intelligenza artificiale per caratterizzare il rumore, ottimizzare i parametri del circuito e migliorare l'affidabilità dei risultati su hardware NISQ.

L'elaborato si sviluppa seguendo un percorso che va dalle fondamenta teoriche del calcolo quantistico fino ai risultati sperimentali su simulatori rumorosi, fornendo al lettore gli strumenti concettuali e pratici necessari per comprendere sia le potenzialità che i limiti dell'algoritmo di Shor nell'era attuale, e le prospettive aperte dall'integrazione con l'intelligenza artificiale. Gli obiettivi specifici e la struttura dettagliata del lavoro sono descritti nel Capitolo~\ref{chap:obiettivi}.

\end{onehalfspace}
